\begin{corollary}
Let $n\ge1$ be an integer, let $1<p<\infty$ and let $K'\ge1$.  Assume the
bundle (C2) of \noderef{QCPackage}.  Let $W\subseteq\mathbb{C}$ be open with
$0\in W$ and let $f,\chi,H\colon\mathbb{C}\to\mathbb{C}$ satisfy
\begin{enumerate}
\item[(H1)] $f$ solves equation (12) on $W$ in the sense of
  \noderef{SolvesEq12}, and $f(0)=0$;
\item[(H2)] $\chi$ is $K'$-quasiconformal on $W$ in the sense of
  \noderef{QuasiconformalOn}, $\chi(0)=0$, and $f(z)=[\chi(z)]^{n}$ for every
  $z\in W$;
\item[(H3)] $\chi$ is $C^\infty$ on $W-\{0\}$;
\item[(H4)] $H(\chi(z))=z$ for every $z\in W$, and $H$ is continuous on
  $\chi(W)$.
\end{enumerate}
Write $\gamma_n$ for the exponent of \noderef{GammaExp} and $D^{w}$ for the
iterated partial derivative of \noderef{PartialWord}.  Then there is $\delta>0$
and there are constants $B_m$, $m=0,1,2,\dots$, such that
\begin{equation}\tag{48}
  \sum_{|w|=m}\bigl|D^{w}f(z)\bigr|\le B_m|z|^{n/\gamma_n-m}
  \qquad\text{for }0<|z|<\delta,
\end{equation}
the sum running over all $2^m$ words $w$ of length $m$ in $\{0,1\}$.

This is Corollary 2 of the paper (paper.tex:473--479).  Hypothesis (H1) is the
paper's ``$f$ is a local solution to the system (12) which has the form (15),
$f(0)=0$'', and (H2) is the factorization (15) itself, whose existence and
whose quasiconformality of $\chi$ are part of the quoted bundle (C2); (H3) and
(H4) are Section 3's normalisation of the hodograph pair.  The words of length
$m$ realize exactly the operators $D^\nu$ with $|\nu|=m$, as
\noderef{PartialWord} records, so the sum is the paper's $\sum_{|\nu|=m}$ up to
the multiplicities with which each $D^\nu$ occurs, at most $2^m$, absorbed into
$B_m$.  Equation (18) for $H$ is not assumed; it is a consequence of (H1) and
(H2).
\end{corollary}
\begin{proof}
Write $\gamma=\gamma_n$ and $V=\chi(W)$.  Note first that $f(0)=0$ is also a
consequence of (H2): $f(0)=\chi(0)^{n}=0$ since $\chi(0)=0$ and $n\ge1$.  It is
recorded in (H1) because the paper states it, and is not used separately below.

\emph{Step 0: $\gamma>0$.}  By part (a) of \noderef{GammaRatioMinAtOne}, with
$n\ge1$ and $1<p$, we have $\gamma>0$.

\emph{Step 1: $V$ is an open neighbourhood of $0$.}  By
\noderef{QuasiconformalOn} applied to $\chi$ on $W$, the set $V=\chi(W)$ is
open, $\chi$ is injective and continuous on $W$.  Since $0\in W$ and
$\chi(0)=0$, we have $0=\chi(0)\in V$.

\emph{Step 2: $H$ is a two-sided inverse of $\chi$ between $V$ and $W$.}  Let
$\xi\in V$, say $\xi=\chi(z)$ with $z\in W$.  By (H4), $H(\xi)=H(\chi(z))=z$,
so
\begin{equation}\tag{1}
  H(\xi)\in W\quad\text{and}\quad \chi(H(\xi))=\chi(z)=\xi
  \qquad(\xi\in V).
\end{equation}
Taking $z=0$ in (H4) and using $\chi(0)=0$ gives
\begin{equation}\tag{2}
  H(0)=H(\chi(0))=0 .
\end{equation}
Also $H(V)=W$: the inclusion $H(V)\subseteq W$ is (1), and conversely each
$z\in W$ is $H(\chi(z))$ with $\chi(z)\in V$.

Finally, $H$ does not vanish on $V-\{0\}$: if $\xi\in V$ and $H(\xi)=0$ then
$\xi=\chi(H(\xi))=\chi(0)=0$ by (1).  Hence
\begin{equation}\tag{3}
  H(\xi)\in W-\{0\}\qquad(\xi\in V-\{0\}).
\end{equation}

\emph{Step 3: the inverse clauses of (C2), and (18) on $V$.}  The clause
\texttt{qc\_inverse} of \noderef{QCPackage}, applied with $K'$, the open set
$W$, the map $\chi$ and the map $H$ --- its hypotheses being that $W$ is open,
that $\chi$ is $K'$-quasiconformal on $W$, which is (H2), and that
$H(\chi(z))=z$ for every $z\in W$, which is (H4) --- gives
\begin{equation}\tag{3$'$}
  H\in W^{1,2}_{\mathrm{loc}}(V),\qquad
  H\ \text{differentiable a.e. on }V,\qquad
  J_H>0\ \text{a.e. on }V,
\end{equation}
where $V=\chi(W)$ and $J_H$ is the Jacobian of \noderef{Jac}; in particular
$J_H\neq0$ almost everywhere on $V$, a positive number being nonzero.

Now the hypotheses of \noderef{HodographEq18} hold with this $n$, $p$, the
open sets $W$ and $V$, the point $z_0=0\in W$, the maps $f,\chi,H$, the
factorization of (H2), $\chi(0)=0$, injectivity of $\chi$ on $W$ from Step 1,
the smoothness of $\chi$ on $W-\{0\}$ from (H3), the second and third clauses
of (3$'$) in place of any smoothness of $H$ --- these are its (H3$'$) --- its
(H3$''$), namely $H\in W^{1,2}_{\mathrm{loc}}(V)$ in the sense of
\noderef{MemWLocOn}, which is the first clause of (3$'$), the two relations
(1), $H(0)=0$ from (2), and equation (12) from (H1).  Therefore
\begin{equation}\tag{18}
  H\ \text{solves (18) on }V\ \text{in the sense of \noderef{SolvesEq18}} .
\end{equation}

\emph{Step 4: $H\in W^{1,2}_{\mathrm{loc}}(V)$.}  This is the first clause of
(3$'$):
\begin{equation}\tag{4}
  H\in W^{1,2}_{\mathrm{loc}}\bigl(\chi(W)\bigr)=W^{1,2}_{\mathrm{loc}}(V)
\end{equation}
in the sense of \noderef{MemWLocOn}.

\emph{Step 5: $H$ is $K$-quasiconformal on $V$ with
$K=\max\{p-1,1/(p-1)\}$.}  Since $p>1$ we have $p-1>0$; if $p\ge2$ then
$K\ge p-1\ge1$, and if $p<2$ then $0<p-1<1$ so $K\ge1/(p-1)>1$.  Thus
$K\ge1$.  We verify the five clauses of \noderef{QuasiconformalOn} for $H$ on
$V$.

\emph{Quasiregularity.}  By \noderef{Eq18ImpliesQuasiregular}, applied with
$p$, $n$, the set $V$ and the map $H$, using (4) and (18), $H$ is
$K$-quasiregular on $V$ with exactly this $K$.

\emph{Injectivity.}  If $\xi_1,\xi_2\in V$ and $H(\xi_1)=H(\xi_2)$ then by (1)
$\xi_1=\chi(H(\xi_1))=\chi(H(\xi_2))=\xi_2$.

\emph{Continuity} on $V=\chi(W)$ is part of (H4).

\emph{Openness of the image.}  $H(V)=W$ by Step 2, which is open.

\emph{The inverse.}  Take $\psi=\chi$: by (1), $\chi(H(\xi))=\xi$ for every
$\xi\in V$, and $\chi$ is continuous on $H(V)=W$ by Step 1.

\emph{Step 6: the derivatives of $\chi$.}  The hypotheses of
\noderef{ChiHigherDerivBound} hold with this $n$, $p$, $K$, the open sets $V$
and $U=W-\{0\}$ --- the latter open since $W$ is --- the point $0\in V$ from
Step 1, the maps $H$ and $\chi$, continuity of $H$ on $V$ from (H4) and its
differentiability almost everywhere on $V$ from (3$'$),
$K$-quasiconformality from Step 5, equation (18) from
Step 3, $H(0)=0$ from (2), smoothness of $\chi$ on $U=W-\{0\}$ from (H3), the
inclusion $H(\xi)\in U$ for $\xi\in V-\{0\}$ from (3), and
$\chi(H(\xi))=\xi$ on $V$ from (1).  It therefore provides $\delta>0$ with
\begin{equation}\tag{5}
  \{z:|z|<\delta\}\subseteq H(V)=W
\end{equation}
and constants $B_l'\ge0$ with
\begin{equation}\tag{6}
  \bigl|D^{w}\chi(z)\bigr|\le B_l'\,|z|^{1/\gamma-l}
  \qquad(|w|=l,\ 0<|z|<\delta).
\end{equation}

\emph{Step 7: the Leibniz rule for words.}  We record the two-factor Leibniz
rule in the form needed.  Let $u,v$ be $C^\infty$ on an open set $G$ and let
$w$ be a word of length $m$, with letters indexed by $\{1,\dots,m\}$.  For
$S\subseteq\{1,\dots,m\}$ let $w_S$ denote the subword of $w$ on the positions
in $S$, in their original order.  Then on $G$
\begin{equation}\tag{7}
  D^{w}(uv)=\sum_{S\subseteq\{1,\dots,m\}}D^{w_S}u\cdot D^{w_{S^{c}}}v .
\end{equation}
This is proved by induction on $m$.  For $m=0$ both sides are $uv$, the only
subset being $\emptyset$.  For the step, write $w=j\!:\!\!:w'$ with $w'$ of
length $m-1$.  By \noderef{PartialWord}, $D^{w}=\partial_j\circ D^{w'}$, so by
the inductive hypothesis and the ordinary product rule for $\partial_j$,
\[
  D^{w}(uv)=\partial_j\Bigl(\sum_{S'}D^{w'_{S'}}u\cdot D^{w'_{S'^{c}}}v\Bigr)
  =\sum_{S'}\Bigl(D^{j::w'_{S'}}u\cdot D^{w'_{S'^{c}}}v
    +D^{w'_{S'}}u\cdot D^{j::w'_{S'^{c}}}v\Bigr),
\]
again using $D^{j::\cdot}=\partial_j\circ D^{\cdot}$.  Subsets $S$ of the
letters of $w$ correspond bijectively to pairs consisting of a subset $S'$ of
the letters of $w'$ together with the decision whether the new letter lies in
$S$; the two summands above are exactly the two cases, since the new letter in
$S$ gives $w_S=j\!:\!\!:w'_{S'}$ and $w_{S^{c}}=w'_{S'^{c}}$, and the new
letter in $S^{c}$ gives the other pair.  This is (7).

\emph{Step 8: the derivatives of $\chi^{N}$.}  We claim that for every integer
$N\ge1$ and every $m\ge0$ there is $A_{N,m}\ge0$ with
\begin{equation}\tag{8}
  \bigl|D^{w}\bigl(\chi^{N}\bigr)(z)\bigr|\le A_{N,m}\,|z|^{N/\gamma-m}
  \qquad(|w|=m,\ 0<|z|<\delta).
\end{equation}
Note first that $\chi$ is $C^\infty$ on the open set
$\{z:0<|z|<\delta\}$: that set is contained in $W-\{0\}$ by (5), where $\chi$
is $C^\infty$ by (H3).  Hence so is every power $\chi^{N}$, and all the
derivatives below exist.

For $N=1$, (8) is (6) with $A_{1,m}=B_m'$.  For the inductive step, apply (7)
on $G=\{z:0<|z|<\delta\}$ with $u=\chi$ and $v=\chi^{N}$, so that
$uv=\chi^{N+1}$, and estimate each term by (6) and (8): for $|w|=m$ and
$0<|z|<\delta$,
\[
  \bigl|D^{w}\bigl(\chi^{N+1}\bigr)(z)\bigr|
  \le\sum_{S}\bigl|D^{w_S}\chi(z)\bigr|\,\bigl|D^{w_{S^{c}}}(\chi^{N})(z)\bigr|
  \le\sum_{S}B_{|S|}'\,|z|^{1/\gamma-|S|}\cdot A_{N,m-|S|}\,|z|^{N/\gamma-(m-|S|)},
\]
where $|w_S|=|S|$ and $|w_{S^{c}}|=m-|S|$.  The powers of the positive real
$|z|$ combine additively, and
\[
  \Bigl(\frac1\gamma-|S|\Bigr)+\Bigl(\frac N\gamma-m+|S|\Bigr)
  =\frac{N+1}{\gamma}-m ,
\]
so the right-hand side equals $A_{N+1,m}|z|^{(N+1)/\gamma-m}$ with
\[
  A_{N+1,m}=\sum_{S\subseteq\{1,\dots,m\}}B_{|S|}'\,A_{N,m-|S|}
  =\sum_{j=0}^{m}\binom{m}{j}B_j'\,A_{N,m-j}\ \ge0,
\]
a finite sum of products of nonnegative numbers.  This proves (8) for $N+1$,
and hence for every $N\ge1$ by induction.

\emph{Step 9: conclusion.}  Let $m\ge0$, let $w$ be a word of length $m$, and
let $0<|z|<\delta$.  By (5), $z\in W$, and $W$ is open, so $f$ agrees with
$\chi^{n}$ on a neighbourhood of $z$ by (H2); hence
$D^{w}f(z)=D^{w}(\chi^{n})(z)$.  By (8) with $N=n$, which is legitimate since
$n\ge1$,
\[
  \bigl|D^{w}f(z)\bigr|\le A_{n,m}\,|z|^{n/\gamma-m}.
\]
Summing over the $2^{m}$ words of length $m$,
\[
  \sum_{|w|=m}\bigl|D^{w}f(z)\bigr|\le 2^{m}A_{n,m}\,|z|^{n/\gamma-m}
  \qquad(0<|z|<\delta),
\]
which is (48) with $B_m=2^{m}A_{n,m}\ge0$, with the $\delta$ of Step 6.  This
is the paper's closing step at paper.tex:510.
\end{proof}
